% Configurazione
\documentclass{article}

\usepackage{titling} % Required for inserting the subtitle
\usepackage{graphicx} % Required for inserting images
\usepackage{tabularx} % Per l'ambiente tabularx (tabelle)
\usepackage{calc} % Sempre per le tabelle
\usepackage[hidelinks]{hyperref} % Per i collegamenti ipertestuali, ad esempio sulla table of contents
\usepackage[italian]{babel} % Per la lingua italiana nelle scritte automatiche
\usepackage{xcolor} % Per colorare il testo
\usepackage{lipsum} % Per generare lorem ipsum
\usepackage[normalem]{ulem} % Per sottolineare il testo
\usepackage{array} % Per la visualizzazione fluttuante di array di domande e risposte
\usepackage{ragged2e} % Pacchetto necessario per \justifying che giustifica il testo di tabelle

\newcommand{\ulhref}[2]{\href{#1}{\uline{#2}}} % Nuovo comando per sottolineare i link
\newcommand{\ulref}[1]{\uline{\ref{#1}}} % Nuovo comando per sottolineare i collegamenti a immagini

\usepackage{eurosym} %inserisce il simbolo dell'euro
\usepackage{tikz} %Per creare grafici e disegni
\usepackage{pgf-pie} %Per creare grafico a torta

\graphicspath{ {immagini/} }

% Variabili
\newcommand{\titolo}{Documento degli Impegni}


% Struttura
\begin{document}

\begin{minipage}{0.4\textwidth}
    \includegraphics[width=0.6\textwidth]{Logo Universita di Padova.png}
\end{minipage}
\begin{minipage}{0.55\textwidth}
    \textcolor{red}{\textbf{Università degli Studi di Padova}} \\
    \textcolor{red}{Laurea: Informatica} \\
    \textcolor{red}{Corso: Ingegneria del Software} \\
    \textcolor{red}{Anno Accademico: 2024/2025}
\end{minipage}

\begin{minipage}{0.4\textwidth}
    \includegraphics[width=0.6\textwidth]{Logo SWEg.png}
\end{minipage}
\begin{minipage}{0.55\textwidth}
    \textbf{Gruppo: SWEg Labs} \\
    Email: \textsf{gruppo.sweg@gmail.com}
\end{minipage}

\vspace{2cm}

\centering\section*{\Huge \titolo}
\vspace{0.5cm}

Con il presente documento il gruppo SWEg intende comunicarLe che si impegner\`a a terminare il progetto didattico entro il giorno DD/MM/YYYY.

Dopo un'attenta discussione siamo riusciti a elaborare un preventivo relativo al capitolato... .\\
Il gruppo ha deciso di dividersi equamente le ore di lavoro per ciascun ruolo da assumere durante svolgimento del progetto.\\
La seguente tabella indica le ore di lavoro dei vari ruoli che si occupano del progetto e dei loro relativi costi.

\vspace{0.5cm}

\begin{table}[h]
\centering
\begin{tabular}{|c|c|c|c|c|}
\hline
Ruolo & Ore Ind. & Ore Totali & Costo (\euro/h) & Costo Totale (\euro) \\
\hline
Responsabile & 9 & 63 & 30 & 1.890 \\
\hline
Amministratore & 8 & 56 & 20 & 1.120 \\
\hline
Analista & 5 & 35 & 25 & 875 \\
\hline
Progettista & 23 & 161 & 25 & 4.025 \\
\hline
Programmatore & 24 & 168 & 15 & 2.520 \\
\hline
Verificatore & 26 & 182 & 15 & 2.730 \\
\hline
\end{tabular}
\caption{Ripartizione oraria e dettaglio dei costi}
\end{table}
\newpage


\end{document}